\chapter*{Production and Evidence Method}
\addcontentsline{toc}{chapter}{Production and Evidence Method}
\index{evidence status}\index{AI-assisted production}\index{revision pin}

This edition was assembled and revised with extensive AI-agent assistance operating against
pinned repositories, source code, tests, project records, and generated verification artifacts.
Agents helped inventory APIs, compare claims across chapters, inspect implementation paths,
draft and edit prose, generate structured diagrams, run validators, and challenge the resulting
manuscript. AI output was not treated as technical authority. The cited repository state,
external specification or implementation, test oracle, and retained artifact determine each
claim.

The edition is deliberately revision-bounded. Source behavior is described at the pins in
Appendix~\ref{ch:repo-map}; executed results additionally name their configuration and host.
Automated source scans, builds, PDF checks, and adversarial review passes reduce drift but cannot
make a changing codebase timeless. Later editions should correct stale statements when the
repositories or their evidence change. Responsibility for the published selection, wording, and
claims remains with the named human author and editor.

\section*{Evidence labels used in this edition}

\begin{center}
\small
\begin{tabularx}{\textwidth}{>{\raggedright\arraybackslash}p{3.0cm}X}
\toprule
\textbf{Label} & \textbf{Meaning} \\
\midrule
Source-proven & A reachable implementation, fallback, refusal, or absence was established from
  the pinned source. No successful build or execution is implied. \\
Compile-proven & Names, signatures, includes, templates, and required link closure were exercised
  in the named configuration. \\
Runtime-observed & The named executable path ran and produced the recorded state, value, file,
  trace, or lifecycle event. \\
Renderer-engaged & The intended renderer family and native or translation route were shown to
  execute, rather than a substitute or early skip path. \\
Oracle-compared & The observation was compared with an identified authority under a stated
  normalization or tolerance policy. \\
Pixel-verified & Selected pixels or a complete frame were compared at the stated readback boundary,
  renderer, host, and tolerance. \\
Manually observed & A human inspected the result, but no retained discriminating oracle supplied
  an automated verdict. \\
Historically recorded & A dated project artifact reports the result; this edition did not rerun it. \\
Blocked / unsupported / not attempted & Three distinct outcomes: an external precondition
  prevented the attempt, the named path has no implementation, or no attempt was made. \\
\bottomrule
\end{tabularx}
\end{center}

These labels are coordinates, not a single quality ladder. A compile probe can answer an API
question that a screenshot cannot; a hostile-input test and a pixel oracle exercise different
contracts. Appendix~\ref{ch:verification-tiers} gives the full evidence vector and preferred
authority for each kind of claim.

Repository paths and revisions serve as the primary references, so the book does not maintain a
separate bibliography database. External links are few and reviewed. Source locations describe
the pin and may move later; names, revision identifiers, and the repository map provide the
durable reproduction route.

\section*{PDF accessibility boundary}

The source supplies a consistent heading hierarchy, document language and metadata, bookmarks,
selectable text, Unicode mappings, and descriptions for figures. The current pdfLaTeX pipeline
does not produce a fully tagged PDF, so this edition does not claim PDF/UA conformance. Moving to
a tagged-PDF toolchain remains publishing work for a later release and requires a separate
structural and assistive-technology review.

